%\ifodd\value{page}\hbox{}\newpage\fi
\chapter*{}

\begin{sloka}
{\sanskritfont
\begin{tabular}{l}
लमित्यादिपञ्चपूजा \\
``लं'' पृथिवीतत्त्वात्मिकायै श्रीललितादेव्यै गन्धं परिकल्पयामि ।\\
``हं'' आकाशतत्त्वात्मिकायै श्रीललितादेव्यै पुष्पं परिकल्पयामि ।\\
``यं'' वायुतत्त्वात्मिकायै श्रीललितादेव्यै धूपं परिकल्पयामि ।\\
``रं'' वह्नितत्त्वात्मिकायै श्रीललितादेव्यै दीपं परिकल्पयामि ।\\
``वं'' अमृततत्त्वात्मिकायै श्रीललितादेव्यै अमृतनैवेद्यं परिकल्पयामि ।\\
``सं'' सर्वतत्त्वात्मिकायै श्रीललितादेव्यै सर्वोपचारान् परिकल्पयामि ।
\end{tabular}
}
\end{sloka}

\vspace{1em}

{\scriptsize \iastfont
lamityādi-pañcapūjā\\
``laṃ'' pṛthivī-tattvātmikāyai śrī-lalitā-devyai gandhaṃ parikalpayāmi |\\
``haṃ'' ākāśa-tattvātmikāyai śrī-lalitā-devyai puṣpaṃ parikalpayāmi |\\
``yaṃ'' vāyu-tattvātmikāyai śrī-lalitā-devyai dhūpaṃ parikalpayāmi |\\
``raṃ'' vahni-tattvātmikāyai śrī-lalitā-devyai dīpaṃ parikalpayāmi |\\
``vaṃ'' amṛta-tattvātmikāyai śrī-lalitā-devyai amṛta-naivedyaṃ parikalpayāmi |\\
``saṃ'' sarva-tattvātmikāyai śrī-lalitā-devyai sarvopacārān parikalpayāmi |
}

\vspace{1em}

{\scriptsize
\anvaya{%
लमित्यादिपञ्चपूजा\\
पृथिवीतत्त्वात्मिकायै श्रीललितादेव्यै “लं” गन्धं परिकल्पयामि ।
आकाशतत्त्वात्मिकायै श्रीललितादेव्यै “हं” पुष्पं परिकल्पयामि ।
वायुतत्त्वात्मिकायै श्रीललितादेव्यै “यं” धूपं परिकल्पयामि ।
वह्नितत्त्वात्मिकायै श्रीललितादेव्यै “रं” दीपं परिकल्पयामि ।
अमृततत्त्वात्मिकायै श्रीललितादेव्यै “वं” अमृतनैवेद्यं परिकल्पयामि ।
सर्वतत्त्वात्मिकायै श्रीललितादेव्यै “सं” सर्वोपचारान् परिकल्पयामि ।
}
}

\wordmeaning{%
\wordline{लमित्यादिपञ्चपूजा}{lamityādi-pañcapūjā}
  { la pūjā (serie di offerte rituali) che inizia con il bīja “laṃ” e così via.}%

\wordline{लं}{laṃ - bīja-mantra dell’elemento terra;}{}%
\wordline{पृथिवी-तत्त्व-आत्मिकायै}{pṛthivī-tattva-ātmikāyai}
  { a colei la cui essenza è il principio della terra;}%
\wordline{गन्धं}{gandhaṃ}
  { profumo;}%

\wordline{हं}{haṃ - bīja-mantra dell’elemento spazio;}{}%
\wordline{आकाश-तत्त्व-आत्मिकायै}{ākāśa-tattva-ātmikāyai}
  { a colei la cui essenza è il principio dello spazio;}%
\wordline{पुष्पं}{puṣpaṃ}
  { fiore;}%

\wordline{यं}{yaṃ - bīja-mantra dell’elemento aria;}{}%
\wordline{वायु-तत्त्व-आत्मिकायै}{vāyu-tattva-ātmikāyai}
  { a colei la cui essenza è il principio dell’aria;}%
\wordline{धूपं}{dhūpaṃ}
  { incenso;}%

\wordline{रं}{raṃ -  bīja-mantra dell’elemento fuoco;}{}%
\wordline{वह्नि-तत्त्व-आत्मिकायै}{vahni-tattva-ātmikāyai}
  { a colei la cui essenza è il principio del fuoco;}%
\wordline{दीपं}{dīpaṃ}
  { lampada;}%

\wordline{वं}{vaṃ - bīja-mantra del principio dell’amṛta;}{}%
\wordline{अमृत-नैवेद्यं}{amṛta-naivedyaṃ}
  { offerta rituale di nettare;}%

\wordline{सं}{saṃ - bīja-mantra della totalità;}{}%
\wordline{सर्व-उपचारान्}{sarva-upacārān}
  { tutti gli atti di culto, tutti gli onori rituali;}%

\wordline{श्री-ललिता-देव्यै}{śrī-lalitā-devyai}
  { alla Dea Śrī Lalitā;}%
\wordline{परिकल्पयामि}{parikalpayāmi}
  { io offro, io presento.}%
}
